\begin{pastexam}[label={px:ncu-107-dsa-09}]{中央資工所 107 學年度 第 9 題}
（單選；每題 4 分，答錯倒扣一分）

Suppose that array A[1:n] maintains a binary tree. For a binary tree node stored in A[i], its two
children (if exist) are stored in A[2i] and A[2i+1], respectively. the program below aims to
complete the following two tasks
\begin{enumerate}[label=(\arabic*)]
\item adjust array A to establish a max heap, and
\item apply heap sort on the max heap built in (1) in nondecreasing order.
\end{enumerate}
\begin{minted}[fontsize=\small,breaklines]{c}
void adjust (int A[], int root, int n)
{
  int child, rootkey;
  int temp;
  temp = A[root];
  rootkey = A[root];
  child = 2* root;
  while ( child <=n ) {
     if (( child <=n ) && ( A[child] < A[child+1] ))
        child ++;
     if ( rootkey > A[child] )
        break;
     else {
        ____(B1)____;
        child *=2; }
  }
    ____(B2)____;
}

void heapsort (int A[], int n)
{ /* perform a heap sort on A[1:n] */
   int i, j;
   int temp;
   for ( ___(B3)___; i>0; i++)  /* adjust the binary tree to establish the max heap */
      adjust (A, i, n);
   for (1=n-1; i>0; i--) {      /* heap sort */
      swap( A[1], A[i+1], temp);  /* exchange A[i] and A[i+1] */
      adjust(A, 1, i); }
}
\end{minted}

Blank (B2) in the program above should be\\
(a)~\texttt{A[child] = temp}\\
(b)~\texttt{A[root] = A[child]}\\
(c)~\texttt{A[child] = A[child/2]}\\
(d)~\texttt{A[child/2] = temp}\\
(e)~\texttt{A[root] = temp}
\end{pastexam}

\begin{pastsoln}
答案是 (d) \texttt{A[child/2] = temp}。

上一題（歷屆試題~\ref{px:ncu-107-dsa-08}）說明了 \texttt{adjust} 是挖洞的寫法：原來的
\texttt{A[root]} 存在 \texttt{temp}，迴圈裡洞的位置一直是 \texttt{child/2}。 迴圈結束有兩種情形：
\begin{itemize}
\item \texttt{break}：\texttt{rootkey > A[child]}，\texttt{temp} 比洞底下較大的子節點還大，就放在洞裡。
\item \texttt{child > n}：洞已經是葉節點，沒有子節點可比，\texttt{temp} 也放在洞裡。
\end{itemize}
兩種情形洞都在 \texttt{child/2}，所以 (B2) 是 \texttt{A[child/2] = temp}。 上一題的追蹤最後
\texttt{child = 8}，\texttt{A[4] = temp = 2}。

逐項檢查：
\begin{itemize}
\item (a) \texttt{A[child] = temp}：\texttt{child} 是洞的子節點，可能已經超過 $n$（上一題的例子是 $8 > 6$），
寫到堆積外面，洞也沒填，錯。
\item (b) \texttt{A[root] = A[child]}：\texttt{temp} 不見了，錯。
\item (c) \texttt{A[child] = A[child/2]}：洞裡的值是已經往上搬走的舊值，往下抄沒有意義，\texttt{temp} 也
不見了，錯。
\item (d) 對。
\item (e) \texttt{A[root] = temp}：只有第一輪就 \texttt{break} 時（洞還在 \texttt{root}）才對；洞往下移過時，
會蓋掉搬上來的最大值，錯。 上一題的例子會變成 $2, 6, 7, 6, 5, 8$。
\end{itemize}
\end{pastsoln}
